%% ============================================================================
%% End matter — Elsevier mandatory closing sections.
%%
%% Order per Elsevier author guidelines:
%%   CRediT author statement
%%   Declaration of Competing Interest
%%   Acknowledgements
%%   Funding
%%   Ethics statement
%%   Data availability
%%   Reproducibility
%% ============================================================================

%% NOTE: CRediT authorship contribution statement is auto-generated by
%% cas-sc's \printcredits (called after this file in main.tex). It reads
%% the per-author \credit{...} macros in the frontmatter. Do NOT duplicate
%% the section here. The taxonomy citation~\citep{Brand2015CRediT} is
%% retained in the acknowledgements paragraph below.

\section*{Declaration of Competing Interest}

The authors declare that they have no known competing financial interests or
personal relationships that could have appeared to influence the work reported
in this paper.

\section*{Acknowledgements}

We thank our colleagues at the Department of Computer Science and Engineering,
Shiv Nadar University, Chennai, for detailed feedback on the pre-registered
design of experiments and on the interim gate-check results. We are indebted
to reviewers of earlier iterations of the framework for feedback that shaped
the human-evaluation protocol and threats-to-validity structure adopted here.
We thank the two independent annotators (DR, RR) who joined the three
author-team annotators in conducting the five-annotator human-evaluation
study for their calibration and rating effort. Author contributions follow
the CRediT taxonomy~\citep{Brand2015CRediT}. This work would not have been
feasible without the open-weight releases from Meta, Microsoft, Alibaba,
Google, Mistral, and DeepSeek, and the maintainers of the Ollama runtime,
the LangChain / LangGraph ecosystem, and the BERTScore, HumanEval, and MBPP
benchmarks.

\section*{Funding}

This research received no specific grant from any funding agency in the public,
commercial, or not-for-profit sectors. The compute cost of the experimental
sweep was covered by the corresponding author's personal Colab Pro+
subscription.

\section*{Ethics statement}

\paragraph{Data.} All datasets used in this study are publicly available prior
work: HumanEval~\citep{Chen2021HumanEval}, MBPP~\citep{Austin2021MBPP}, and
LiveCodeBench~\citep{Jain2025LiveCodeBench}. No new function-level data was
collected. The reference test suites, docstrings, and code accompanying these
benchmarks are used as-is.

\paragraph{Human participants.} The 60-pair human-evaluation study described in
Section~\ref{sec:experiment_setup_human_eval} and reported in
Section~\ref{sec:results_human_eval} involved five human annotators rating pairs
of code artifacts (docstrings, code fragments, and test suites): three from the
author team (unpaid; contribution recorded via CRediT), and two external
annotators recruited independently of the pipeline runners and compensated at
a rate consistent with regional software-engineering hourly rates. The
distinction between author-team and external annotators is analytically
material: five-way inter-rater $\alpha_{\textsf{ord}} = 0.193$ is lower than the
three-way (author-team-only) $\alpha_{\textsf{ord}} = 0.376$, a decline
attributable to shared calibration among author-team members
(Section~\ref{sec:human_eval_agreement_results}). No personally-identifying
data, biometric data, health data, or demographic data was collected from any
annotator. The annotation protocol was reviewed under the institutional research
ethics framework at the corresponding author's institution; the review
determined that the study was exempt from full IRB review under the
``research involving benign educational tests or observations of public
behavior'' category. Annotator identities are not linked to their ratings in
any published artifact.

\paragraph{AI systems.} All experiments were conducted with open-weight
language models. No closed-weight models were used in the primary results. LLM
outputs were treated strictly as artifacts under evaluation --- the LLM systems
themselves are not the subject of ethical claims. The judge SLM (DeepSeek-R1:14B)
was used exclusively as an evaluation tool, not as a decision-maker for any
human-consequential decision.

\paragraph{Environmental impact.} The experimental sweep consumed approximately
22~GPU-days on Google Colab hardware ($\sim$530~GPU-hours total), split roughly
40\% on NVIDIA A100 (initial mono sweep and pilot) and 60\% on NVIDIA T4
(majority of hetero and null cells). Assuming sustained-load power draws of
$\sim$350~W (A100) and $\sim$70~W (T4), a data-center PUE of $1.10$, and the
Google Cloud global-average grid intensity of $\sim$240~gCO$_2$-eq/kWh, we
estimate total emissions of approximately \textbf{25~kg~CO$_2$-equivalent},
with an uncertainty range of $\sim$20--50~kg CO$_2$-equivalent depending on
regional grid mix. We disclose this as a matter of methodological transparency
and to enable readers to weigh the environmental cost against the study's
scientific yield. All numbers were computed following the MLCO$_2$
estimator~\citep{Lacoste2019MLCO2Emissions} methodology.

\section*{Data availability}

The complete replication package for this study is publicly available at
\url{https://github.com/balajivenky06/roundtrip-closure} under the MIT
license. The repository contains: the full experimental sweep TSV
(\texttt{results/results\_roundtrip.tsv}, 5{,}890 measurement rows covering
all 20 cells of the pre-registered DOE); the 30-function pilot TSV; the
pre-registered DOE table (\texttt{doe.py}); the model lineup
(\texttt{config.py}); the three closure-path drivers (\texttt{closure\_paths.py});
the closure-validity decision (\texttt{closure\_decision.py}, Algorithm~2); the
test-filter validity gate (\texttt{closure\_metrics.filter\_tests\_with\_reason},
Algorithm~3); the mutation-testing engine, following standard mutation-analysis methodology~\citep{Andrews2005MutationAnalysis, Just2014MutationAnalysis}
(\texttt{mutation\_testing.py}); the 67-test suite covering the closure-validity
decision, the test-filter validity gate, and their integration into the sweep
pipeline (\texttt{tests/test\_closure\_decision.py},
\texttt{tests/test\_filter\_tests\_reason.py}, \texttt{tests/test\_wiring\_integration.py});
the analysis pipeline (\texttt{scripts/run\_analysis.py}); all six gap-closer
scripts (\texttt{scripts/build\_*.py}); the 24 LaTeX tables and 5 figures
(1~TikZ framework diagram + 4~empirical plots) used in this manuscript, all
reproducible from the sweep TSV via a single command; the 60-pair
five-annotator human-evaluation TSVs (\texttt{human\_eval/data/annotator\_*.tsv})
with the pair-metadata table; and the three-judge frontier-judge replay outputs
(\texttt{human\_eval/data/frontier\_judge\_*.tsv}) from GPT-4o-mini, Claude
Haiku~4.5, and Llama-3.3-70B.

The pre-processed HumanEval, MBPP, LiveCodeBench, and HumanEval-Mutated datasets
are included in the \texttt{data/} directory as JSONL files. The five-layer
checkpointing infrastructure (LLM call cache, pytest subprocess cache,
BERTScore embedding cache, TSV row cache, JSONL decontamination cache)
preserves every intermediate artifact and enables exact re-execution of any
single path on any single sample without re-invoking the underlying SLMs.

\section*{Code and reproducibility}

The single command
\begin{verbatim}
python3 scripts/run_analysis.py --tsv results/results_roundtrip.tsv
\end{verbatim}
regenerates the 24 LaTeX tables and 5 figures from the sweep TSV. The
following commands regenerate the six gap-closer artifacts described in
Section~5:

\begin{itemize}
  \item G1 (threshold sensitivity):
        \texttt{python3 scripts/build\_threshold\_sensitivity.py}
  \item G2 (per-mutation-operator, requires LLM cache):
        \texttt{python3 scripts/regen\_per\_operator.py} followed by
        \texttt{python3 scripts/build\_per\_operator\_table.py}
  \item G3 (judge-metric disagreement):
        \texttt{python3 scripts/build\_disagreement\_decomposition.py}
  \item G4 (capability preservation):
        \texttt{python3 scripts/build\_capability\_preservation.py}
  \item G5 (path~$\times$~cell interaction ANOVA):
        \texttt{python3 scripts/build\_path\_cell\_anova.py}
  \item G6 (stage-contribution figure):
        \texttt{python3 scripts/build\_stage\_contribution\_figure.py}
\end{itemize}

Re-executing the full 20-cell experimental sweep from cold is possible via
\texttt{python3 scripts/run\_pilot.py} (30-function pilot) followed by
\texttt{python3 train\_roundtrip.py} per cell. Full-sweep wall-clock cost on
Colab Pro+ A100: approximately 15--22~hours; on T4: 60--90~hours. The
five-layer checkpointing means partial re-execution is inexpensive.

\paragraph{Cache-efficiency accounting.} The full 20-cell sweep consumed
16{,}302 SLM calls, of which 3{,}706 were served from cache ($22.7\%$ hit rate).
Total wall-clock time on Google Colab hardware was approximately 22~GPU-days
over 23~calendar days, spanning both A100 and T4. Mono cells run first and
see near-zero hit rates ($\leq 0.6\%$ for M1, M2, M4, M5, M6); M3 achieves
$21.5\%$ after a mid-sweep re-run. Hetero cells inherit substantial hit rates:
H4 achieves $89.2\%$ because its qwen-coder stages reuse M6 cache entries.
N2 and N3 achieve $100\%$ because their surviving stages match earlier cells.
Full per-cell hit rates in \texttt{tables/tab\_cache\_efficiency.tex}.

%% Preprint disclosure section omitted: no preprint has been posted.
%% If a preprint is posted later (e.g. arXiv), reinstate this block with
%% the preprint URL, posting date, and a list of substantive changes.
